% Ad Familiares 5.20 (ad-familiares-05-20)
% Source: https://raw.githubusercontent.com/PerseusDL/canonical-latinLit/master/data/phi0474/phi056/phi0474.phi056.perseus-lat1.xml
% Fetched by scripts/fetch_latin.py

% Dateline (Perseus): Scr. ad urbem paulo post prid. Non. Ian. a. 705 (49).

\ciceroLetterOpener{CICERO RVFO.}

\ciceroSection{1} quoquo modo potuissem, te convenissem, si eo, quo constitueras, venire voluisses. qua re, etsi mei commodi causa commovere me noluisti, tamen ita existimes velim, me antelaturum fuisse, si ad me misisses, voluntatem tuam commodo meo. ad ea, quae scripsisti, commodius equidem possem de singulis ad te rebus scribere, si M. Tullius, scriba meus, adesset. A quo mihi exploratum est in rationibus dumtaxat referendis (de ceteris rebus adfirmare non possum) nihil eum fecisse scientem, quod esset contra aut rem aut existimationem tuam; dein, si rationum referendarum ius vetas et mos antiquus maneret, me relaturum rationes, nisi tecum pro coniunctione nostrae necessitudinis contulissem confecissemque, non fuisse.

\ciceroSection{2} quod igitur fecissem ad urbem, si consuetudo pristina maneret, id, quoniam lege Iulia relinquere rationes in provincia necesse erat easdemque totidem verbis referre ad aerarium, feci in provincia, neque ita feci ut te ad meum arbitrium adducerem, sed tribui tibi tantum, quantum me tribuisse numquam me paenitebit. totum enim scribam meum, quem tibi video nunc esse suspectum, tibi tradidi. tu ei M. Mindium, fratrem tuum, adiunxisti. rationes confectae me absente sunt tecum; ad quas ego nihil adhibui praeter lectionem. ita accepi librum a meo servo scriba, ut eundem acceperim a fratre tuo. si honos is fuit, maiorem tibi habere non potui, si fides, maiorem tibi habui quam paene ipsi mihi; si providendum fuit ne quid aliter ac tibi et honestum et utile esset referretur, non habui cui potius id negoti darem, quam cui dedi. illud quidem certe factum est, quod lex iubebat, ut apud duas civitates, Laudicensem et Apamensem, quae nobis maxime videbantur, quoniam ita necesse erat, rationes confectas conlatas deponeremus. itaque huic loco primum respondeo me, quamquam iustis de causis rationes deferre properarim, tamen te exspectaturum fuisse, nisi in provincia relictas rationes pro relatis haberem, quam ob rem

\ciceroSection{3} de Volusio quod scribis, non est id rationum; docuerunt enim me periti homines, in his cum omnium peritissimus tum mihi amicissimus, C. Camillus, ad Volusium traferri nomen a Valerio non potuisse, praedes Valerianos teneri. neque id erat HS x_x_x_ , ut scribis, sed HS x_i_x_ ; erat enim curata nobis pecunia Valeri mancipis nomine; ex qua reliquum quod erat in rationibus rettuli.

\ciceroSection{4} sed sic me et liberalitatis fructu privas et diligentiae et, quod minime tamen laboro, mediocris etiam prudentiae: liberalitatis, quod mavis scribae mei beneficio quam meo legatum meum praefectumque Q. Leptam maxima calamitate levatos, cum praesertim non deberent esse obligati; diligentiae, quod existimas de tanto officio meo, tanto etiam periculo, nec scisse me quicquam nec cogitavisse, scribam, quicquid voluisset, cum id mihi ne recitavisset quidem, rettulisse; prudentiae, †cum rem a me non insipienter excogitatam quidem putas. nam et Volusi liberandi meum fuit consilium, et, ut multa tam gravis Valerianis praedibus ipsique T. Mario depelleretur, a me inita ratio est; quam quidem omnes non solum probant, sed etiam laudant, et, si verum scire vis, hoc uni scribae meo intellexi non nimium placere. sed ego putavi esse viri boni, cum populus suum servaret, consulere fortunis tot vel amicorum vel civium.

\ciceroSection{5} nam de Lucceio est ita actum, ut auctore Cn. Pompeio ista pecunia in fano poneretur. id ego agnovi meo iussu esse factum. qua pecunia Pompeius est usus ut ea, quam tu deposueras, Sestius. sed haec ad te nihil intellego pertinere; illud me non animadvertisse moleste ferrem, ut ascriberem te in fano pecuniam iussu meo deposuisse, nisi ista pecunia gravissimis esset certissimisque monumentis testata, cui data, quo senatus consulto, quibus tuis, quibus meis litteris P. Sestio tradita esset. quae cum viderem tot vestigiis impressa, ut in iis errari non posset, non ascripsi id, quod tua nihil referebat; ego tamen ascripsisse mallem, quoniam id te video desiderare.

\ciceroSection{6} sicut scribis tibi id esse referendum, idem ipse sentio, neque in eo quicquam a meis rationibus discrepabunt tuae; addes enim tu, meo iussu, quod ego, qui non addidi, nec causa est cur negem nec, si causa esset et tu nolles, negarem. nam de sestertiis nongentis milibus certe ita relatum est, ut tu sive frater tuus referri voluit. sed , si quid est, quoniam de Lucceio parum provisum est, quod ego in rationibus referendis etiam nunc corrigere possim, de eo mihi, quoniam senatus consulto non sum usus, quid per leges liceat, considerandum est. te certe in pecuniam exactam ista referre ex meis rationibus relatis non oportuit, nisi quid me fallit; sunt enim alii peritiores. illud cave dubites, quin ego omnia faciam, quae interesse tua aut etiam velle te existimem, si ullo modo facere possim.

\ciceroSection{7} quod scribis de beneficiis scito a me et tribunos militaris et praefectos et contubernalis dumtaxat meos delatos esse. in quo quidem me ratio fefellit; liberum enim mihi tempus ad eos deferendos existimabam dari, postea certior sum factus triginta diebus deferri necesse esse, quibus rationes rettulissem. sane moleste tuli non illa beneficia tuae potius ambitioni reservata esse quam meae, qui ambitione nihil uterer. de centurionibus tamen et de tribunorum militarium contubernalibus res est in integro; genus enim horum beneficiorum definitum lege non erat.

\ciceroSection{8} reliquum est de sestertiis centum milibus, de quibus memini mihi a te Myrina litteras esse adlatas, non mei errati, sed tui; in quo peccatum videbatur esse, si modo erat, fratris tui et Tulli. sed , cum id corrigi non posset, quod iam depositis rationibus ex provincia decessimus, credo me quidem tibi pro animi mei voluntate proque ea spe facultatum, quam tum habebamus, quam humanissime potuerim, rescripsisse; sed neque tum me humanitate litterarum mearum obligatum puto neque me tuam hodie epistulam de HS c_ sic accepisse, ut ii accipiunt, quibus epistulae per haec tempora molestae sunt.

\ciceroSection{9} simul illud cogitare debes, me omnem pecuniam, quae ad me salvis legibus pervenisset, Ephesi apud publicanos deposuisse; id fuisse HS x_x_i_i_ ; eam omnem pecuniam Pompeium abstulisse. quod ego sive aequo animo sive iniquo fero, tu de HS c_ aequo animo ferre debes et existimare eo minus ad te vel de tuis cibariis vel de mea liberalitate pervenisse. quod si mihi expensa ista HS c tulisses, tamen, quae tua est suavitas quique in me amor, nolles a me hoc tempore aestimationem accipere; nam, numeratum si cuperem, non erat. sed haec iocatum me putato, ut ego te existimo. ego tamen, cum Tullius rure redierit, mittam eum ad te, si quid ad rem putabis pertinere. hanc epistulam cur conscindi velim causa nulla est.
